% GENERATED FILE. Edit canonical lessons, then run npm run generate:books.
% canonical-source-hash: fnv1a64:6abcb9e4bf440e0c
% canonical-lessons: UR-C325-dev, UR-C325-pari, UR-C325-jadu, UR-C325-bhut, UR-C325-dain

\chapter{Giants and Fairies}
\label{ch:ur-giants-and-fairies}
\hlchaptermodality{325}

\begin{chapteropening}
\textbf{By the end of this chapter:} \emph{I can tell a story with magic in it.}

Complete the last lesson of chapter 325: I can tell a story with magic in it.
\end{chapteropening}

\section[\texorpdfstring{\ur{دیو} (dev)}{دیو (dev)}]{\ur{دیو} (dev) — a giant in tales; a demon}

\label{lesson:UR-C325-dev}



\begin{quote}
Before the new one: say the Urdu for a riddle, then the Urdu for a proverb.
\end{quote}

\subsection*{You'll want to know: \ur{دیو}}
\textbf{\ur{دیو}} — \emph{dev} — \textquotedblleft{}a giant in tales; a demon\textquotedblright{}.

\begin{quote}
\ur{پہاڑ} \ur{پر} \ur{ایک} \ur{دیو} \ur{رہتا} \ur{تھا۔}
\end{quote}

\begin{quote}
\emph{pahāṛ par ek dev rahtā thā.}
\end{quote}

\begin{quote}
A giant lived on the mountain.
\end{quote}

Masculine.

\begin{cousinweb}
is usually traced to Persian \emph{dēv}, \textquotedblleft{}a demon\textquotedblright{}, the word inside \textbf{\ur{دیوانہ}}. In Urdu tales the \textbf{\ur{دیو}} is a huge, cruel giant.
\end{cousinweb}

\begin{grammarlens}[title={What you've built}]
One more word for telling what happened.
\end{grammarlens}

\subsection*{Your turn}
\begin{itemize}
\raggedright
  \item \emph{Say it:} \emph{dev}
  \item \emph{Say it:} \emph{dev}, once more
  \item \emph{Say it:} \emph{pahāṛ par ek dev rahtā thā.} — A giant lived on the mountain.
  \item \emph{Recall:} say the Urdu for a riddle, then the Urdu for a proverb, then say \emph{dev} again
\end{itemize}

\begin{tcolorbox}[breakable,colback=teal!4,colframe=teal!35!black,title={Before you move on}]
What does \ur{دیو} mean? (\textquotedblleft{}A giant in tales; a demon\textquotedblright{}.) Say it once more.
\end{tcolorbox}
\section[\texorpdfstring{\ur{پری} (parī)}{پری (parī)}]{\ur{پری} (parī) — a fairy}

\label{lesson:UR-C325-pari}



\begin{quote}
Before the new one: say the Urdu for a proverb, then the Urdu for a giant in tales.
\end{quote}

\subsection*{You'll want to know: \ur{پری}}
\textbf{\ur{پری}} — \emph{parī} — \textquotedblleft{}a fairy\textquotedblright{}.

\begin{quote}
\ur{پری} \ur{نے} \ur{بچے} \ur{کو} \ur{ایک} \ur{تحفہ} \ur{دیا۔}
\end{quote}

\begin{quote}
\emph{parī ne bachche ko ek tohfā diyā.}
\end{quote}

\begin{quote}
The fairy gave the child a gift.
\end{quote}

Feminine.

\begin{cousinweb}
\textbf{\ur{پری}} (\emph{parī}) is usually traced to Persian \emph{parī}, \textquotedblleft{}a fairy\textquotedblright{}. Urdu calls a beautiful face \textbf{\ur{پری} \ur{جیسا}}.
\end{cousinweb}

\begin{grammarlens}[title={What you've built}]
One more word for telling what happened.
\end{grammarlens}

\subsection*{Your turn}
\begin{itemize}
\raggedright
  \item \emph{Say it:} \emph{parī}
  \item \emph{Say it:} \emph{parī}, once more
  \item \emph{Say it:} \emph{parī ne bachche ko ek tohfā diyā.} — The fairy gave the child a gift.
  \item \emph{Recall:} say the Urdu for a proverb, then the Urdu for a giant in tales, then say \emph{parī} again
\end{itemize}

\begin{tcolorbox}[breakable,colback=teal!4,colframe=teal!35!black,title={Before you move on}]
What does \ur{پری} mean? (\textquotedblleft{}A fairy\textquotedblright{}.) Say it once more.
\end{tcolorbox}
\section[\texorpdfstring{\ur{جادو} (jādū)}{جادو (jādū)}]{\ur{جادو} (jādū) — magic}

\label{lesson:UR-C325-jadu}



\begin{quote}
Before the new one: say the Urdu for a giant in tales, then the Urdu for a fairy.
\end{quote}

\subsection*{You'll want to know: \ur{جادو}}
\textbf{\ur{جادو}} — \emph{jādū} — \textquotedblleft{}magic\textquotedblright{}.

\begin{quote}
\ur{اس} \ur{نے} \ur{جادو} \ur{سے} \ur{پتھر} \ur{کو} \ur{پھول} \ur{بنایا۔}
\end{quote}

\begin{quote}
\emph{us ne jādū se patthar ko phūl banāyā.}
\end{quote}

\begin{quote}
She turned the stone into a flower by magic.
\end{quote}

Masculine.

\begin{cousinweb}
\textbf{\ur{جادو}} (\emph{jādū}) is usually traced to Persian \emph{jādū}, \textquotedblleft{}magic, a sorcerer\textquotedblright{}. \textbf{\ur{جادوگر}} is a magician.
\end{cousinweb}

\begin{grammarlens}[title={What you've built}]
One more word for telling what happened.
\end{grammarlens}

\subsection*{Your turn}
\begin{itemize}
\raggedright
  \item \emph{Say it:} \emph{jādū}
  \item \emph{Say it:} \emph{jādū}, once more
  \item \emph{Say it:} \emph{us ne jādū se patthar ko phūl banāyā.} — She turned the stone into a flower by magic.
  \item \emph{Recall:} say the Urdu for a giant in tales, then the Urdu for a fairy, then say \emph{jādū} again
\end{itemize}

\begin{tcolorbox}[breakable,colback=teal!4,colframe=teal!35!black,title={Before you move on}]
What does \ur{جادو} mean? (\textquotedblleft{}Magic\textquotedblright{}.) Say it once more.
\end{tcolorbox}
\section[\texorpdfstring{\ur{بھوت} (bhūt)}{بھوت (bhūt)}]{\ur{بھوت} (bhūt) — a ghost}

\label{lesson:UR-C325-bhut}



\begin{quote}
Before the new one: say the Urdu for a fairy, then the Urdu for magic.
\end{quote}

\subsection*{You'll want to know: \ur{بھوت}}
\textbf{\ur{بھوت}} — \emph{bhūt} — \textquotedblleft{}a ghost\textquotedblright{}.

\begin{quote}
\ur{بچوں} \ur{کو} \ur{بھوت} \ur{سے} \ur{ڈر} \ur{لگتا} \ur{ہے۔}
\end{quote}

\begin{quote}
\emph{bachchoṅ ko bhūt se ḍar lagtā hai.}
\end{quote}

\begin{quote}
Children are afraid of ghosts.
\end{quote}

Masculine.

\begin{cousinweb}
\textbf{\ur{بھوت}} (\emph{bhūt}) is inherited from Sanskrit \emph{bhūta}, \textquotedblleft{}a being, a spirit\textquotedblright{} (Monier-Williams): what has been, and haunts.
\end{cousinweb}

\begin{grammarlens}[title={What you've built}]
One more word for telling what happened.
\end{grammarlens}

\subsection*{Your turn}
\begin{itemize}
\raggedright
  \item \emph{Say it:} \emph{bhūt}
  \item \emph{Say it:} \emph{bhūt}, once more
  \item \emph{Say it:} \emph{bachchoṅ ko bhūt se ḍar lagtā hai.} — Children are afraid of ghosts.
  \item \emph{Recall:} say the Urdu for a fairy, then the Urdu for magic, then say \emph{bhūt} again
\end{itemize}

\begin{tcolorbox}[breakable,colback=teal!4,colframe=teal!35!black,title={Before you move on}]
What does \ur{بھوت} mean? (\textquotedblleft{}A ghost\textquotedblright{}.) Say it once more.
\end{tcolorbox}
\section[\texorpdfstring{\ur{ڈائن} (ḍāin)}{ڈائن (ḍāin)}]{\ur{ڈائن} (ḍāin) — a witch}

\label{lesson:UR-C325-dain}



\begin{quote}
Before the new one: say the Urdu for magic, then the Urdu for a ghost.
\end{quote}

\subsection*{You'll want to know: \ur{ڈائن}}
\textbf{\ur{ڈائن}} — \emph{ḍāin} — \textquotedblleft{}a witch\textquotedblright{}.

\begin{quote}
\ur{ڈائن} \ur{جنگل} \ur{میں} \ur{رہتی} \ur{تھی۔}
\end{quote}

\begin{quote}
\emph{ḍāin jangal meṅ rahtī thī.}
\end{quote}

\begin{quote}
The witch lived in the forest.
\end{quote}

Feminine.

\begin{cousinweb}
\textbf{\ur{ڈائن}} (\emph{ḍāin}) is inherited: CDIAL lists the Hindi cousin \emph{ḍāin}, \textquotedblleft{}a witch\textquotedblright{}, under Sanskrit \emph{ḍākínī} (CDIAL 5542).
\end{cousinweb}

\begin{grammarlens}[title={What you've built}]
One more word for telling what happened.
\end{grammarlens}

\subsection*{Your turn}
\begin{itemize}
\raggedright
  \item \emph{Say it:} \emph{ḍāin}
  \item \emph{Say it:} \emph{ḍāin}, once more
  \item \emph{Say it:} \emph{ḍāin jangal meṅ rahtī thī.} — The witch lived in the forest.
  \item \emph{Recall:} say the Urdu for magic, then the Urdu for a ghost, then say \emph{ḍāin} again
\end{itemize}

\begin{tcolorbox}[breakable,colback=teal!4,colframe=teal!35!black,title={Before you move on}]
What does \ur{ڈائن} mean? (\textquotedblleft{}A witch\textquotedblright{}.) Say it once more.
\end{tcolorbox}
